\section{Vanishing lower mass}
\label{sec:vanishing-lower-mass}

By \eqref{eq:gaud-lower-mass-formula}, the lower mass $v_-(h)$ vanishes as soon as there are Gauduchon metrics $\eta$ whose degrees $\int_Xh\wedge\eta^{n-1}$ stay bounded while their volumes $\int_X\eta^n$ tend to infinity. We construct such metrics on a compact threefold. Writing $\Gamma=\eta^2$, the volume of $\eta$ is a determinant integral of the coefficient matrix of the positive $\ddc$-closed form $\Gamma$. The construction first concentrates the trace of a normalized form $\Gamma$ in one direction near a hypersurface, and then adds an average of compact rational curves that fills the two complementary directions. The degree stays bounded, while the determinant, and hence the volume of the positive root, becomes large. The smooth equation then converts these metrics into potentials whose masses tend to zero.

\begin{theorem}\label{thm:vanishing-lower-mass}\index{Positive lower mass!failure}
There is a connected compact complex threefold $X$ such that, for every Hermitian form $h$ on $X$, there are smooth real functions $u_j$ satisfying
\[
 h+\ddc u_j>0,\qquad \sup_Xu_j=0,\qquad
 \int_X(h+\ddc u_j)^3\longrightarrow0.
\]
In particular, $v_-(h)=0$ for every Hermitian form $h$ on this manifold.
\end{theorem}

\subsection{The elliptic quotient}\label{subsec:lower-mass-elliptic-quotient}

Write points of $\mathbb P^3$ as $[z:w]$, with $z,w\in\C^2$, and set
\[
 \Omega=\{[z:w]:z\ne0,\ w\ne0\},\qquad
 q=\mathrm{e}^{2\pi},\qquad g[z:w]=[z:qw].
\]
Thus $g$ is an automorphism of $\Omega$; let $\langle g\rangle$ be the cyclic group it generates. This group acts freely. Indeed, if $g^k$ fixed a point, the same scalar would have to equal $1$ on the nonzero block $z$ and $q^k$ on the nonzero block $w$, which forces $k=0$.

\Needspace{8\baselineskip}
The function
\[
 t=\log\frac{\|w\|^2}{\|z\|^2}
\]
is well defined on $\Omega$ and satisfies $t\circ g=t+4\pi$. Consequently the action is properly discontinuous: $t$ is bounded on compact subsets, so only finitely many translates of one compact set can meet another. Moreover, the strip $\{0\leq t\leq4\pi\}$ is compact. In homogeneous coordinates of total norm one, its defining inequalities bound both blocks away from zero, so it is a closed subset of $\mathbb P^3$ contained in $\Omega$. Since the strip meets every orbit,
\[
 X=\Omega/\langle g\rangle
\]
is a connected compact complex threefold.

The map $[z:w]\mapsto([z],[w])$ induces a holomorphic submersion
$X\to\mathbb P^1\times\mathbb P^1$. Using a logarithm of the scalar relating two points with the same image, one sees that its fibers are the elliptic curves
$E=\C/(2\pi\mathbb Z+2\pi\mathrm{i}\mathbb Z)$. Let $\alpha,\beta$ be the pullbacks of the Fubini--Study forms of the two factors, each normalized to have integral one. Although $t$ does not descend to $X$, its differentials $\mathrm dt$ and $\dc t$ do. The form
\begin{equation}\label{eq:lower-mass-quotient-metric}
 \omega=\alpha+\beta+\mathrm{d}t\wedge\dc t
\end{equation}
is Hermitian: $\alpha+\beta$ is semipositive with kernel the vertical line, while $\mathrm{d}t\wedge\dc t$ is semipositive and strictly positive on the vertical line.

\Needspace{7\baselineskip}
On the open set $\{z_0w_0\ne0\}$ we use the coordinates
\[
 x=\frac{z_1}{z_0},\qquad y=\frac{w_1}{w_0},\qquad
 \zeta=\log\frac{w_0}{z_0}\in E.
\]
Near $0\in E$, the last coordinate is represented by an ordinary complex number. For any one of these coordinates, denoted by $v$, put
$\sigma_v=\frac{\mathrm{i}}{2\pi}\mathrm{d}v\wedge\mathrm{d}\bar v$.
Choose a smooth probability density $p$ relative to $\sigma_v$ that is positive on $\{|v|<1\}$ and vanishes on $\{|v|\geq1\}$. Choose also a smooth function $\chi_0$ with values in $[0,1]$, supported in $\{|v|<1\}$ and equal to one on $\{|v|\leq1/2\}$, and write $\chi(x,y)=\chi_0(x)\chi_0(y)$. For $0<\delta<1/100$ define
\[
 B_\delta=\chi(x,y)\delta^{-2}p(\zeta/\delta)\sigma_\zeta,
\]
extended by zero. This is a smooth semipositive $(1,1)$-form on $X$, because the chosen functions are smooth with compact support in their coordinate domains. As $\delta\downarrow0$ it concentrates on the hypersurface $\{\zeta=0\}$.

The first step shows that normalized positive $\ddc$-closed $(2,2)$-forms can pair arbitrarily strongly with $B_\delta$.

\begin{lemma}\label{lem:lower-mass-trace-concentration}
For $\delta>0$ put
\[
 C_\delta=\sup\left\{\int_XB_\delta\wedge\Theta:
 \Theta>0,\ \ddc\Theta=0,\ \int_X\omega\wedge\Theta=1\right\},
\]
the supremum being taken over smooth real $(2,2)$-forms $\Theta$. These numbers are finite and tend to infinity as $\delta\downarrow0$.
\end{lemma}
\Needspace{12\baselineskip}
\begin{proof}
For fixed $\delta$, choose a constant $M_\delta>0$ with $B_\delta\leq M_\delta\omega$; this shows that $C_\delta\leq M_\delta$. Suppose instead that $C_{\delta_l}\leq C$ along a sequence $\delta_l\to0$. Then $\int_X(C\omega-B_{\delta_l})\wedge\Theta\geq0$ for every admissible $\Theta$, and in particular for the squares of Gauduchon metrics. By \cref{thm:lamari}, there are quasi-plurisubharmonic functions $v_l$ such that
\[
 C\omega+\ddc v_l\geq B_{\delta_l},\qquad \sup_Xv_l=0.
\]
By \cref{thm:normalized-compactness}, a subsequence converges in $L^1$ to a $C\omega$-plurisubharmonic function $v$, not identically $-\infty$. Since $\delta^{-2}p(\zeta/\delta)\sigma_\zeta$ is an approximate identity in the variable $\zeta$,
\[
 B_{\delta_l}\rightharpoonup\chi[\zeta=0],\qquad
 C\omega+\ddc v\geq\chi[\zeta=0].
\]
On a coordinate ball where $\chi=1$, add to $v$ a smooth strictly plurisubharmonic function whose complex Hessian dominates $C\omega$. The resulting plurisubharmonic function has a current that dominates $[\zeta=0]$, so its Lelong number is at least one at every point of an open subset of this hypersurface. Smooth corrections do not change Lelong numbers, so the same holds for $v$. Siu's analyticity theorem, recalled in \cref{subsec:psh-lelong-levels}, shows that
$Z=\{\nu(v,\cdot)\geq1\}$ is an analytic subset of $X$. It is a proper subset because $v$ is not identically $-\infty$.

Let $\widetilde Z$ be the inverse image of $Z$ in $\Omega$. It contains an open subset of
$H_0\cap\Omega$, where $H_0=\{w_0=z_0\}$. This surface is a projective plane with two points removed, so it is connected and irreducible, and the identity principle gives $H_0\cap\Omega\subset\widetilde Z$. Since $\widetilde Z$ is invariant under $g$, also
\[
 H_k\cap\Omega\subset\widetilde Z,\qquad H_k=\{w_0=q^kz_0\},
 \qquad k\in\mathbb Z.
\]
Now use the chart $z_1=1$ near the point $[0:1:1:0]\in\Omega$. There $H_k$ is the graph $z_0=q^{-k}w_0$. For each fixed $(w_0,w_1)$ near $(1,0)$, every local holomorphic defining function of $\widetilde Z$ vanishes at the distinct values $z_0=q^{-k}w_0$, which tend to zero as $k\to+\infty$. The one-variable identity theorem makes each defining function vanish identically on a polydisk. Thus $\widetilde Z$ contains a nonempty open subset of the connected manifold $\Omega$, contradicting properness. Hence no subsequence of $C_\delta$ stays bounded as $\delta\downarrow0$.
\end{proof}

By the lemma we can choose $\delta_j\downarrow0$, each smaller than half its predecessor, and then, by the definition of the supremum, smooth strictly positive $\ddc$-closed $(2,2)$-forms $\Theta_j$ with
\begin{equation}\label{eq:lower-mass-large-trace}
 \int_X\omega\wedge\Theta_j=1,\qquad
 J_j=\int_XB_{\delta_j}\wedge\Theta_j>j^4.
\end{equation}

\subsection{Averaging the trace and compact curves}

A large pairing with $B_{\delta_j}$ means that $\Theta_j$ has a large coefficient in one direction, but possibly only on a small set. Averaging by translations spreads this coefficient over a whole product region of size comparable to $\delta_j$ in the vertical direction. A second average, this time of compact rational curves, fills the two complementary directions, which is needed for a lower bound on a determinant. Fix $j$ and write $\delta=\delta_j$.

For complex parameters $a,b,\tau$, the automorphisms
\[
 T_{a,b,\tau}[z_0:z_1:w_0:w_1]
 =[z_0:z_1+az_0:\mathrm{e}^{\tau}w_0:
                  \mathrm{e}^{\tau}(w_1+bw_0)]
\]
of $\Omega$ commute with $g$, and in the product chart they act by
$(x,y,\zeta)\mapsto(x+a,y+b,\zeta+\tau)$.
\Needspace{13\baselineskip}
Average $T_{a,b,\tau}^*\Theta_j$ against the probability density
\[
 3^{-4}p(a/3)p(b/3)(3\delta)^{-2}p(\tau/(3\delta))
       \sigma_a\wedge\sigma_b\wedge\sigma_\tau,
\]
and denote the result by $\widehat\Theta_j$. It is smooth, strictly positive and $\ddc$-closed. The parameters range in a fixed compact set of automorphisms, so $(T_{a,b,\tau}^{-1})^*\omega\leq C\omega$ with one constant $C$, and a change of variables gives
\begin{equation}\label{eq:lower-mass-averaged-degree}
 \int_X\omega\wedge\widehat\Theta_j\leq C.
\end{equation}

Let $F_j$ be the coefficient of $\sigma_x\wedge\sigma_y$ in the restriction of $\Theta_j$ to the slices $\{\zeta=\text{constant}\}$, and let $\widehat F_j$ be the corresponding coefficient of $\widehat\Theta_j$. Positivity gives $F_j\geq0$, and since the translations preserve $\sigma_x\wedge\sigma_y$, $\widehat F_j$ is an average of translates of $F_j$. At a point of
\[
 V_j=\{|x|,|y|<1/4,\ |\zeta|<\delta/4\},
\]
every translation carrying it into $\operatorname{supp}B_\delta$ has
$|a|,|b|\leq5/4$ and $|\tau|\leq5\delta/4$. On this range the first two averaging factors are bounded below by a positive constant, and the third by a positive constant times $\delta^{-2}$. On the other hand, $\chi\leq1$ and $p$ is bounded, so the number $J_j$ of \eqref{eq:lower-mass-large-trace} is at most a constant times $\delta^{-2}$ times the integral of $F_j$ over $\operatorname{supp}B_\delta$. Comparing the two integrals gives
\begin{equation}\label{eq:lower-mass-spread-trace}
 \widehat F_j\geq cJ_j\quad\hbox{on }V_j,
\end{equation}
with $c>0$ independent of $j$.

The second average fills the two horizontal directions. Put
\[
 f(a,b)=p(4(a-1))p(b/2),\qquad
 \rho_j(\tau)=4\delta^{-2}p(2\tau/\delta).
\]
For parameters in the supports of these functions and $|u|\leq1/16$, consider the holomorphic maps $\iota_{\tau,u,a,b}\colon\mathbb P^1\to X$ defined by
\begin{equation}\label{eq:lower-mass-rational-curves}
 \iota_{\tau,u,a,b}[z_0:z_1]
 =[z_0:z_1:\mathrm{e}^{\tau}(z_0+uz_1):
                \mathrm{e}^{\tau}(az_1+bz_0)]\in X.
\end{equation}
For $0<\varepsilon\leq1/16$, average their integration currents against
\[
 \rho_j(\tau)f(a,b)\varepsilon^{-2}p(u/\varepsilon)
       \sigma_\tau\wedge\sigma_a\wedge\sigma_b\wedge\sigma_u,
\]
and call the result $T_{j,\varepsilon}$. Define $T_{j,0}$ in the same way, but with $u=0$ and without the integral over $u$.

\Needspace{22\baselineskip}
\begin{lemma}\label{lem:lower-mass-curve-average}
For $\varepsilon>0$, the form $T_{j,\varepsilon}$ is smooth on $X$, semipositive and closed, and
\[
 \int_X\omega\wedge T_{j,\varepsilon}\leq C'
\]
with $C'$ independent of $j$ and $\varepsilon$. For each fixed $j$, $T_{j,\varepsilon}$ converges smoothly to $T_{j,0}$ on compact subsets of the product chart as $\varepsilon\downarrow0$. In this chart
\begin{equation}\label{eq:lower-mass-curve-local}
 T_{j,0}=\rho_j(\zeta)\sigma_\zeta\wedge K_f,\qquad
 K_f=\int_\C f(a,y-ax)\frac{\mathrm{i}}{2\pi}
       (\mathrm{d}y-a\,\mathrm{d}x)\wedge
       (\mathrm{d}\bar y-\bar a\,\mathrm{d}\bar x)\,\sigma_a,
\end{equation}
and on $\{|x|,|y|\leq1/4\}$ the form $K_f$ is bounded below by a positive multiple of $\sigma_x+\sigma_y$.
\end{lemma}
\Needspace{8\baselineskip}
\begin{proof}
On the parameter supports, $|a-bu|\geq3/4-2/16=5/8$. Hence the matrix
$\left(\begin{smallmatrix}1&u\\ b&a\end{smallmatrix}\right)$ is uniformly invertible, and the projective line \eqref{eq:lower-mass-rational-curves} lies in $\Omega$. Distinct deck translates of this line are disjoint: equality at a point fixes $[z_0:z_1]$, and invertibility of the matrix then forces the deck scalar to be one. Thus each map is an embedding onto a compact rational curve in $X$.

The evaluation map from $\mathbb P^1$ times the open parameter domain to $X$ is a holomorphic submersion. Indeed, after fixing a projective gauge, variations of $u$ and $a$ span the two $w$-coordinates at points with $z_1\ne0$, variations of $\tau$ and $b$ span them at points with $z_1=0$, and the coordinate on $\mathbb P^1$ supplies the remaining direction. For $\varepsilon>0$, the average is the integral along the fibers of this submersion of a smooth compactly supported parameter density, which proves smoothness. As an average of closed positive integration currents, $T_{j,\varepsilon}$ is closed and semipositive.

All curve parameters lie in one compact subset of the domain where the matrix is invertible, so the $\omega$-areas of the curves are uniformly bounded. The weights in $\tau$ and $u$ have mass one, and the weight in $(a,b)$ has a fixed finite mass. This proves the degree bound.

When $u=0$, the curve is given in the product chart by $\zeta=\tau$ and $y=ax+b$. Integrating in $b$ and $\tau$ yields \eqref{eq:lower-mass-curve-local}; the $(1,1)$-factor there is the integration current of the graph of $y=ax+b$. On $\{|x|,|y|\leq1/4\}$, the density $f(a,y-ax)$ is bounded below by a positive constant for $a$ in a fixed small disk about $1$. The covectors $\mathrm{d}y-a\,\mathrm{d}x$, for $a$ in that disk, have no common nonzero annihilator, so integrating their squared absolute values and using compactness of the unit sphere in the $(x,y)$-directions gives the asserted lower bound.

For general $u$, the curve is given by
\[
 \zeta=\tau+\log(1+ux),\qquad y=\frac{ax+b}{1+ux}.
\]
On a fixed compact subset of the chart, $1+ux$ stays away from zero for all sufficiently small $u$. Solving for the parameters gives
$\tau=\zeta-\log(1+ux)$ and $b=(1+ux)y-ax$. Regarding $\rho_j$ as a smooth function on $E$, this change of variables also applies near the boundary of its support. The Jacobians, the density coefficients and any fixed number of their derivatives depend smoothly on $u$. Averaging against the approximate identity in $u$ therefore gives local $C^\infty$ convergence. This convergence is needed only for each fixed $j$; it is not uniform in $j$.
\end{proof}

\subsection{The volume of the positive root}

The two averages contribute complementary blocks of the coefficient matrix of a $(2,2)$-form: the translated forms give a large vertical entry, and the curves give a positive horizontal block. A Schur-complement estimate combines them into a lower bound for the volume of the positive root. In the coframe $(x,y,\zeta)$ put
$\mathrm{d}V_0=\sigma_x\wedge\sigma_y\wedge\sigma_\zeta$. For a real $(1,1)$-form $\xi$, let $H_\xi$ be its coefficient matrix in this normalization. A real $(2,2)$-form $\Theta$ corresponds to a Hermitian matrix $Q_\Theta$, which we call its cofactor matrix, by
\[
 \xi\wedge\Theta=\operatorname{tr}(H_\xi Q_\Theta)\,\mathrm{d}V_0.
\]
Strict positivity of $\Theta$ is equivalent to $Q_\Theta>0$.

\Needspace{8\baselineskip}
If $\eta^2=\Theta$ and $H=H_\eta$, expanding the wedge products gives
\[
 Q_\Theta=2(\det H)H^{-1},\qquad
 \eta^3=\frac{3}{\sqrt2}(\det Q_\Theta)^{1/2}\,\mathrm{d}V_0.
\]
Conversely, these formulas uniquely recover a smooth positive $H$ from $Q_\Theta>0$; this is the three-dimensional case of the root construction preceding the proof of \cref{thm:lamari}. We write $\Theta^{3/2}=\eta^3$.

\begin{lemma}\label{lem:lower-mass-unbounded-root}
There are smooth strictly positive $\ddc$-closed $(2,2)$-forms $\Gamma_j$ such that
\[
 \int_X\omega\wedge\Gamma_j\leq C_2,\qquad
 R_j=\int_X\Gamma_j^{3/2}\longrightarrow+\infty.
\]
\end{lemma}
\begin{proof}
Write the cofactor matrix of $\widehat\Theta_j$ in block form, with a $2\times2$ Hermitian block $A_j$ in the $(x,y)$-directions and a column $\kappa_j\in\C^2$:
\[
 Q_j=\begin{pmatrix}A_j&\kappa_j\\\kappa_j^*&\widehat F_j\end{pmatrix}.
\]
The cofactor matrix of $T_{j,0}$ has the form $\operatorname{diag}(B_j,0)$. Let $I_2$ denote the $2\times2$ identity matrix. The lower bound for $K_f$ and the definition of $\rho_j$ give a constant $c_0>0$ such that $B_j\geq c_0\delta_j^{-2}I_2$ on $V_j$. Since $Q_j>0$, the Schur complement $A_j-\kappa_j\kappa_j^*/\widehat F_j$ is positive definite, and
\[
 \det\bigl(Q_j+\operatorname{diag}(B_j,0)\bigr)
 =\widehat F_j\det\left(B_j+A_j-
             \frac{\kappa_j\kappa_j^*}{\widehat F_j}\right)
 \geq\widehat F_j\det B_j,
\]
by monotonicity of the determinant on positive Hermitian matrices. Combining the root formula with \eqref{eq:lower-mass-spread-trace}, we obtain
\[
 (\widehat\Theta_j+T_{j,0})^{3/2}
       \geq c\delta_j^{-2}\sqrt{J_j}\,\mathrm{d}V_0
       \quad\hbox{on }V_j.
\]
\Needspace{8\baselineskip}
The smaller region $W_j=\{|x|,|y|<1/8,\ |\zeta|<\delta_j/8\}$ has $\mathrm{d}V_0$-volume equal to a fixed positive constant times $\delta_j^2$. Consequently
\[
 I_j=\int_{W_j}(\widehat\Theta_j+T_{j,0})^{3/2}
       \geq c_1\sqrt{J_j}>c_1j^2.
\]
For this fixed $j$, \cref{lem:lower-mass-curve-average} gives smooth convergence of $T_{j,\varepsilon}$ to $T_{j,0}$ on the closure of $W_j$. The cofactor matrix of $\widehat\Theta_j+T_{j,0}$ is positive definite there, so the integrand of the root volume depends continuously on these smooth coefficients. We can therefore choose $0<\varepsilon_j\leq1/16$ so small that
\[
 \int_{W_j}(\widehat\Theta_j+T_{j,\varepsilon_j})^{3/2}
       \geq I_j/2.
\]
Set $\Gamma_j=\widehat\Theta_j+T_{j,\varepsilon_j}$. It is smooth, strictly positive and $\ddc$-closed. Its degree is at most $C_2=C+C'$, by \eqref{eq:lower-mass-averaged-degree} and \cref{lem:lower-mass-curve-average}, whereas its root volume is at least $c_1j^2/2$.
\end{proof}

\Needspace{15\baselineskip}
\begin{proof}[of \cref{thm:vanishing-lower-mass}]
Fix a Hermitian form $h$ on $X$, and choose $A>0$ with $h\leq A\omega$. Let $\eta_j$ be the positive root of $\Gamma_j$, so that $\eta_j^2=\Gamma_j$. These metrics are Gauduchon, and
\[
 P_j=\int_Xh\wedge\eta_j^2\leq AC_2,\qquad
 R_j=\int_X\eta_j^3\longrightarrow+\infty.
\]
Set $\mu_j=\eta_j^3/R_j$ and solve the smooth equation in the background $h$:
\[
 h_j=h+\ddc u_j>0,\qquad h_j^3=c_h(\mu_j)\mu_j,
 \qquad \sup_Xu_j=0.
\]
Since $D_{\mu_j}(\eta_j)=R_j^{2/3}$, using $\eta_j$ as a test in \cref{thm:gaud-volume-dual-formula} gives
\[
 0<\int_Xh_j^3=c_h(\mu_j)=\mathcal V_{\mu_j}(h)
       \leq\frac{P_j^3}{R_j^2}\longrightarrow0.
\]
The potentials $u_j$ are smooth, hence bounded, so $v_-(h)=0$.
\end{proof}

\subsection{Consequences for dimension and equations}

The example occurs in the smallest possible dimension: on a surface, comparison with a Gauduchon metric gives positive lower mass. Products with projective spaces carry the threefold degeneration to every higher dimension, and a product with the Hopf threefold produces, in addition, unbounded upper mass.

\begin{corollary}\label{cor:lower-mass-dimension}
Every Hermitian form on a compact complex manifold of dimension at most two has positive lower mass. In every dimension $n\geq3$ there is a connected compact complex manifold on which every Hermitian form has zero lower mass.
\end{corollary}
\Needspace{10\baselineskip}
\begin{proof}
In dimension one the mass is $\int_Xh>0$, and the surface case is \cref{cor:geometric-determinant-bound}. In higher dimensions we use the monotonicity and the homogeneity of degree $n$ of $v_-$, which follow directly from \eqref{eq:gaud-lower-mass-formula}, to pass from a split metric to every Hermitian metric.

\Needspace{11\baselineskip}
For $n=3$ this is \cref{thm:vanishing-lower-mass}. For $n>3$, take the product of the threefold $X$ with $\mathbb P^{n-3}$, let $\kappa$ be a Fubini--Study form on the projective factor, and consider the split metric $h_0=\omega+\kappa$, with pullbacks understood. For the potentials $u_j$ given by the theorem for the background $\omega$,
\[
 \int_{X\times\mathbb P^{n-3}}(h_0+\ddc u_j)^n
 =\binom n3\left(\int_X(\omega+\ddc u_j)^3\right)
                   \left(\int_{\mathbb P^{n-3}}\kappa^{n-3}\right)
 \longrightarrow0.
\]
For another Hermitian form $h$, choose $A$ with $h\leq Ah_0$. Monotonicity and homogeneity give $v_-(h)\leq A^n v_-(h_0)=0$.
\end{proof}

\Needspace{16\baselineskip}
\begin{corollary}\label{cor:both-mass-properties-fail}
In every dimension $n\geq6$ there is a connected compact complex manifold with
\[
 v_-(h)=0,\qquad v_+(h)=+\infty
\]
for every Hermitian form $h$.
\end{corollary}
\begin{proof}
Take the product of the threefold above with the Hopf threefold of \cref{chap:hopf}, with a split metric. Potentials pulled back from the first factor make the masses tend to zero, while the smooth potentials of \cref{thm:hopf-unbounded}, pulled back from the second factor, make them tend to infinity. In each case the product integral is $\binom63$ times the product of the two factor masses. The metric comparison in the proof of \cref{cor:lower-mass-dimension} handles the lower mass for every Hermitian form, and the metric independence in \cref{prop:bmp-metric} handles the upper mass. Products with projective spaces give every higher dimension by the same expansion of wedge powers.
\end{proof}

Finally, fix a smooth positive probability measure $\mu$ on the threefold $X$. The smooth equation still gives $\mathcal V_\mu(h)>0$ for each Hermitian form $h$. With $\mathscr M$ and $\lambda_*$ as in \cref{sec:gaud-volume-extrema}, \cref{thm:gaud-volume-mass-extrema} gives
\[
 \inf_{\nu\in\mathscr M}\mathcal V_\nu(h)=0,
 \qquad \lambda_*(h)=0.
\]
Thus even a smooth strictly positive background can have vanishing infimum over densities. For the metrics $\eta_j$ and numbers $R_j,P_j$ in the proof of \cref{thm:vanishing-lower-mass}, the measures $\mu_j=\eta_j^3/R_j$ satisfy
$\mathcal V_{\mu_j}(h)\leq P_j^3/R_j^2\to0$.
Hence the densities $\mu_j/\mu$ cannot be uniformly bounded above: a bound $\mu_j\leq M\mu$ would give $\mathcal V_{\mu_j}(h)\geq M^{-1}\mathcal V_\mu(h)>0$, by the comparison of denominators in the proof of \cref{prop:gaud-volume-density-comparison}.

The example also shows that the endpoint estimate in \cref{thm:gaud-determinant-endpoint} fails for the full family of normalized Gauduchon currents: the supremum of the $L^{3/2}$ norms of their determinant densities is infinite, although every subcritical $L^q$ norm, $1<q<3/2$, remains uniformly bounded by \cref{thm:gaud-determinant-subcritical}. For the same fixed background, \cref{thm:gaud-local-adjoint-endpoint,cor:gaud-local-density-family} give uniform $L^\infty$ bounds for the exact maximizers and their adjoint probabilities in a small $C^{2,\alpha}$ family of smooth densities. A larger sufficient family is given by \cref{thm:gaud-hessian-adjoint-endpoint}: uniformly comparable smooth solution forms yield uniform bounds for both densities, without derivative bounds on the varying Hessian. Endpoint estimates for arbitrary bounded families of continuous densities remain open.
